  % ======================================================================
  %  LiDAR perception pipeline --- technical report
  %  Built from the personal template (template/doc.tex).
  %  Compile: pdflatex -> bibtex -> pdflatex x2   (or: latexmk -pdf technical_report.tex)
  % ======================================================================

  \documentclass[10pt,a4paper]{scrartcl}

  % --------------------------------------------------
  % Encoding and language
  % --------------------------------------------------
  \usepackage[T1]{fontenc}
  \usepackage[utf8]{inputenc}
  \usepackage[english]{babel}

  % --------------------------------------------------
  % Page layout
  % --------------------------------------------------
  \usepackage[
    top=18mm,
    bottom=18mm,
    left=20mm,
    right=20mm,
    includefoot
  ]{geometry}

  % --------------------------------------------------
  % Mathematics and graphics
  % --------------------------------------------------
  \usepackage{graphicx}
  \usepackage{amsmath,amssymb}
  \usepackage{bm}
  \setlength{\jot}{0.3cm}

  % --------------------------------------------------
  % Tables
  % --------------------------------------------------
  \usepackage{booktabs}
  \usepackage{tabularx}
  \usepackage{array}
  \usepackage{ragged2e}
  \usepackage{float}
  \newcolumntype{L}[1]{>{\RaggedRight\arraybackslash}p{#1}}

  % --------------------------------------------------
  % Typography
  % --------------------------------------------------
  \usepackage[scaled=0.95]{helvet}
  \usepackage[expansion=false]{microtype}
  \renewcommand{\familydefault}{\sfdefault}

  % --------------------------------------------------
  % Lists
  % --------------------------------------------------
  \usepackage{enumitem}
  \setlist[itemize]{leftmargin=*, topsep=2pt, parsep=0pt, partopsep=0pt, itemsep=3pt}
  \setlist[enumerate]{leftmargin=*, topsep=2pt, parsep=0pt, partopsep=0pt, itemsep=3pt}
  \setlist[description]{leftmargin=0pt, topsep=2pt, parsep=0pt, itemsep=3pt}

  % --------------------------------------------------
  % Colour, figures, boxes and hyperlinks
  % --------------------------------------------------
  \usepackage[dvipsnames]{xcolor}
  \usepackage{tikz}
  \usetikzlibrary{arrows.meta,positioning,calc}
  \usepackage{pgfplots}
  \pgfplotsset{compat=1.18}
  \usepackage{tcolorbox}
  \tcbuselibrary{skins,breakable}
  \usepackage{caption}
  \usepackage[hidelinks]{hyperref}

  \definecolor{accent}{RGB}{0,129,88}
  \definecolor{leftc}{RGB}{37,99,235}
  \definecolor{rightc}{RGB}{202,138,4}
  \renewcommand{\emph}[1]{\textcolor{accent}{#1}}
  \captionsetup{font=small,labelfont={bf,color=accent}}

  \newtcolorbox{keybox}[1][Key results]{enhanced, breakable,
    colback=accent!5, colframe=accent, boxrule=0.6pt, arc=2pt,
    left=6pt, right=6pt, top=3pt, bottom=3pt,
    fonttitle=\bfseries\small, title={#1}}
  \newtcolorbox{warnbox}[1][Limitation]{enhanced, breakable,
    colback=BrickRed!4, colframe=BrickRed, boxrule=0.6pt, arc=2pt,
    left=6pt, right=6pt, top=3pt, bottom=3pt,
    fonttitle=\bfseries\small, title={#1}}

  \newcommand{\code}[1]{\texttt{#1}}
  \newcommand{\file}[1]{\texttt{\small #1}}

  % --------------------------------------------------
  % Citations
  % --------------------------------------------------
  \usepackage[numbers,square,sort&compress]{natbib}

  % --------------------------------------------------
  % Header and footer
  % --------------------------------------------------
  \usepackage[automark,headsepline=false,footsepline=false]{scrlayer-scrpage}
  \clearpairofpagestyles
  \ifoot{Victor Gabriele De Matteis}
  \cfoot{\pagemark}
  \ofoot{LiDAR Perception --- Technical Report}
  \setlength{\footskip}{20mm}
  \pagestyle{scrheadings}

  % --------------------------------------------------
  % Spacing and sections
  % --------------------------------------------------
  \setlength{\parindent}{0pt}
  \setlength{\parskip}{0.45em}

  \RedeclareSectionCommand[beforeskip=1.4ex plus 0.2ex minus 0.1ex, afterskip=0.6ex]{section}
  \RedeclareSectionCommand[beforeskip=1.0ex plus 0.2ex minus 0.1ex, afterskip=0.4ex]{subsection}
  \RedeclareSectionCommand[beforeskip=0.7ex plus 0.2ex minus 0.1ex, afterskip=0.3ex]{subsubsection}
  \setkomafont{section}{\large\bfseries\color{accent}}
  \setkomafont{subsection}{\normalsize\bfseries\color{accent}}
  \setkomafont{subsubsection}{\normalsize\bfseries}

  \newcommand{\RuleSection}[1]{%
    \section{#1}%
    \vspace{0.2em}{\color{accent!35}\hrule height 0.4pt}\vspace{0.3em}%
  }

  % ======================================================================
  \begin{document}
  % ======================================================================

  \begin{center}
    {\Huge\bfseries\color{accent} LiDAR Perception for Driverless Racing}\\[0.25cm]
    {\large Technical report: cone detection, track boundaries, racing line and odometry}\\[0.2cm]
    {\small E-Agle Formula Student Driverless perception task (LiDAR path) --- October 2026}\\[0.15cm]
    {\small Victor Gabriele De Matteis}\\[0.2cm]
    {\color{accent!35}\hrule height 0.5pt}
  \end{center}

  \vspace{0.3em}

  \begin{abstract}
  \noindent
  This report describes a LiDAR-only perception pipeline for a Formula Student Driverless
  car, written in C++20 with PCL, Eigen and NLopt, and what it achieves. From a single
  24,000-point scan the pipeline finds the track cones, assigns them to the left or right
  boundary, and computes a minimum-curvature racing line. A second program estimates the
  car's motion between two scans with ICP. A Qt viewer shows every stage and lets an operator
  tune the parameters live. Because the recorded data has no labels, the work is validated
  against \emph{ground truth that is known by construction}: synthetic tracks, synthetic
  motions applied to a real scan, and a ray-casting LiDAR simulator. Four automated tests run
  under \code{ctest}, and two evaluation tools write visual reports for human review. The
  validation found and fixed three defects that gave plausible-looking but wrong results. It
  also measured where the system stands: track-boundary and racing-line logic are correct
  on every clean test track, and odometry recovers known motions to 2\,mm. The weak point is
  cone detection, with 37\,\% recall on simulated scans. The full pipeline needs 0.49\,s per
  scan on one CPU core, against a 0.1\,s real-time budget.
  \end{abstract}

  % ======================================================================
  \RuleSection{Summary}
  \label{sec:summary}

  \begin{keybox}[Key results (measured on an Apple M3, Release build)]
  \small
  \begin{tabularx}{\linewidth}{@{}L{6.2cm} X@{}}
  Recorded scan (\code{cones.pcd}) & 24,000 points $\to$ 67 clusters $\to$ 6 cones: 2 left, 3 right, 1 left unassigned \\
  Left/right split, synthetic tracks & 10/10 clean tracks correct (previous version: 2/10) \\
  Racing line, known optima & centreline on a straight (error 0); outer edge on an arc (within 2.3\,\% of the half-width) \\
  Odometry, known motions & error 2\,mm and $<0.01^\circ$ (previous settings: 0.77--0.96\,m, $\approx10^\circ$) \\
  Cone detection, simulated scans & precision 79\,\%, recall 37\,\%, centroid error 6.9\,cm \\
  Runtime, cone pipeline & $485\pm22$\,ms per scan (2.1\,Hz), one core \\
  Runtime, odometry ICP & 2.20\,s per scan pair (median of 10) \\
  Automated tests & 4 \code{ctest} targets, all passing; 2 report generators \\
  \end{tabularx}
  \end{keybox}

  \textbf{What works.} The pipeline runs end to end on the recorded scan and produces a
  usable result. The geometry downstream of detection (boundary assignment, racing line,
  frame conventions, odometry transform) is verified against exact answers.

  \textbf{What does not, yet.} Cone recall falls quickly with range, and the cone/non-cone
  decision does not separate cones from cone-sized objects. Section~\ref{sec:detection}
  quantifies both, together with a measured alternative. Throughput is about 5$\times$ below a
  10\,Hz sensor rate.

  \textbf{Maturity.} A research prototype. It is validated in simulation and on one recorded
  scan, has regression tests for its geometric core, and has a clear, measured list of next
  steps (Section~\ref{sec:roadmap}).

  % ======================================================================
  \RuleSection{Scope and data}
  \label{sec:scope}

  The task (\file{docs/task\_brief.md}) asks for five levels on either camera or LiDAR data:
  load and preprocess, detect cones, classify cones versus obstacles, extract the track or
  racing line, and estimate odometry, plus an interactive debugging GUI. The LiDAR path was
  chosen. Levels 1--4 use \code{cones.pcd}; level 5 uses \code{first.pcd} and \code{second.pcd}.

  Each scan holds 24,000 points. All $x$ values are non-negative, so the sensor is
  forward-looking, and it sits about 0.82\,m above the ground. The data contains no labels,
  no cone colours (intensity is not used) and no timestamps. After the loading transform,
  every module uses one frame: \textbf{$x$ forward, $y$ up, $z$ right}; the track lies in the
  $(x,z)$ plane. Every module was made to use this one convention; an earlier version mixed
  two conventions (Section~\ref{sec:fixes}).

  % ======================================================================
  \RuleSection{System architecture}
  \label{sec:arch}

  \begin{figure}[htbp]
  \centering
  \begin{tikzpicture}[
      node distance=3.2mm,
      box/.style={draw=accent, fill=accent!6, rounded corners=2pt, align=left,
                  text width=7.3cm, inner sep=3.5pt, font=\scriptsize},
      arr/.style={-{Stealth[length=2mm]}, accent, thick}]
    \node[box] (a1) {\textbf{Load + level} (fixed rotation to $x$ fwd, $y$ up, $z$ right)\hfill\textbf{24,000 pts}};
    \node[box, below=of a1] (a2) {\textbf{Height crop} $y\in[-10,-0.2]$ m\hfill\textbf{9,089}};
    \node[box, below=of a2] (a3) {\textbf{Statistical outlier removal} ($k=50$, $\alpha=1.1$)\hfill\textbf{8,251}};
    \node[box, below=of a3] (a4) {\textbf{Region growing} on normals: ground and walls removed\hfill\textbf{4,025}};
    \node[box, below=of a4] (a5) {\textbf{Euclidean clustering} (5 cm)\hfill\textbf{67 clusters}};
    \node[box, right=6mm of a1] (b1) {\textbf{Geometric gates} (size, ground contact, height)\hfill\textbf{9 candidates}};
    \node[box, below=of b1] (b2) {\textbf{ICP against a cone model}, fitness + yaw gate\hfill\textbf{6 cones}};
    \node[box, below=of b2] (b3) {\textbf{Left/right split} (side-consistent greedy growth)\hfill\textbf{2 + 3}};
    \node[box, below=of b3] (b4) {\textbf{Racing line} (bounded min-curvature, BOBYQA)\hfill\textbf{2 gates}};
    \node[box, below=of b4] (b5) {\textbf{Obstacles}: scan minus cones, filtered\hfill\textbf{5,731 pts}};
    \draw[arr] (a1) -- (a2); \draw[arr] (a2) -- (a3); \draw[arr] (a3) -- (a4); \draw[arr] (a4) -- (a5);
    \draw[arr] (a5.east) -- ++(3mm,0) |- (b1.west);
    \draw[arr] (b1) -- (b2); \draw[arr] (b2) -- (b3); \draw[arr] (b3) -- (b4); \draw[arr] (b4) -- (b5);
  \end{tikzpicture}
  \caption{Cone pipeline on \code{cones.pcd} with the GUI default parameters, with the
  measured count after each stage. Odometry is a separate program
  (Section~\ref{sec:odometry}).}
  \label{fig:pipeline}
  \end{figure}

  The repository builds two programs. \file{GUIAndRacingLine} is the Qt viewer that runs the
  cone pipeline; its logic lives in free functions under \file{logic/}, one file per stage.
  \file{OdometryWithProcessing} is a command-line program around a small odometry library
  (\file{odometry.h}). Tests and evaluation tools are in each project's \file{tests/} folder.

  % ======================================================================
  \RuleSection{Methods}
  \label{sec:methods}

  \begin{description}
    \item[] \textbf{Preprocessing.} A fixed rotation levels the scan. A height pass-through
      filter and statistical outlier removal (SOR) follow \citep{rusu2008towards}.
    \item[] \textbf{Ground and wall removal.} Region growing on surface normals
      \citep{rabbani2006segmentation} finds smooth regions of at least 500 points. Cones never
      reach that size (the largest has 113 points), so every such region is removed.
    \item[] \textbf{Clustering.} Euclidean cluster extraction with a 5\,cm tolerance and 8--500
      points per cluster \citep{rusu2009thesis}.
    \item[] \textbf{Cone recognition.} Clusters must have at least 10 points, touch the ground and
      be 0.2--0.5\,m tall. Each survivor is aligned to a synthetic cone surface with
      point-to-point ICP \citep{besl1992icp} and accepted if the mean squared residual is below
      0.005\,m$^2$ and the rotation about the vertical axis is below 13$^\circ$.
    \item[] \textbf{Left/right split.} Seeds: the closest cone on each side of the car. Each
      boundary then grows by its nearest remaining cone that lies on the correct side of the
      \emph{other} boundary and closer to its own, with the turn limited to $\pm45^\circ$. The
      two side conditions are mutually exclusive, so a boundary cannot jump across the track.
    \item[] \textbf{Racing line.} Left and right cones are paired into gates. The decision
      variables are the lateral offsets $d_i$ along each gate's normal, bounded by the half-width
      $w_i$. The cost is
      \begin{equation}
        J(\bm d)=\sum_i \kappa_i^2+0.1\sum_i (d_{i+1}-d_i)^2,\qquad
        \kappa_i=\frac{2\,\theta_i}{\lVert\bm a_i\rVert+\lVert\bm b_i\rVert},\qquad |d_i|\le w_i ,
      \end{equation}
      where $\theta_i=\operatorname{atan2}(\bm a_i\times\bm b_i,\ \bm a_i\cdot\bm b_i)$ is the
      turning angle between consecutive segments $\bm a_i,\bm b_i$. The car is a fixed first
      point, so the line leaves the car smoothly. The problem is solved with the derivative-free,
      bound-constrained BOBYQA \citep{powell2009bobyqa} from NLopt \citep{johnson_nlopt}.
    \item[] \textbf{Odometry.} Both scans are cropped, levelled and SOR-filtered, then registered
      with point-to-point ICP (70 iterations, correspondences limited to 1\,m). PCL returns
      $T_{21}$, which maps the first scan into the second. The car's displacement in the first
      frame is the translation of $T_{21}^{-1}$, i.e.\ $-R^\top\bm t$.
  \end{description}

  % ======================================================================
  \RuleSection{Validation approach}
  \label{sec:validation}

  The recorded scans have no labels, so accuracy cannot be measured on them. Instead, every
  module is tested on inputs whose correct output is \emph{known by construction}:

  \begin{table}[htbp]
  \small
  \centering
  \renewcommand{\arraystretch}{1.1}
  \begin{tabularx}{\textwidth}{@{}L{3.3cm} L{5.3cm} X@{}}
  \toprule
  \textbf{Module} & \textbf{Ground truth} & \textbf{Tool} \\
  \midrule
  Left/right split & 10 synthetic tracks (straights, curves, S-curve, hairpin, car offset and yawed), clean and with noise or missing cones, 100 random seeds each & \code{eval\_left\_right} (ctest + HTML report) \\[.15cm]
  Racing line & Analytic optima (straight, constant arc), exact interpolation geometry, bounds & \code{test\_racing\_line} (ctest) \\[.15cm]
  Cone ICP & Known rotations of the cone model; simulated single views of cones and cone-sized non-cones & \code{test\_cone\_icp} (ctest) \\[.15cm]
  Odometry & A known motion applied to \code{first.pcd}, with independent sampling, 1\,cm noise and partial overlap & \code{test\_odometry} (ctest) \\[.15cm]
  Whole cone pipeline & Ray-cast LiDAR scans (24,000 rays) of 5 scenes with cones, walls, obstacles and cone-sized decoys & \code{eval\_pipeline} (HTML report) \\
  \bottomrule
  \end{tabularx}
  \caption{Validation matrix.}
  \label{tab:validation}
  \end{table}

  The two HTML reports draw every case from above: true and predicted sides, wrong cones
  circled, missed cones and false positives marked. A human reviewer can then confirm the
  metrics, and check the recorded scan, for which no metric exists.

  % ======================================================================
  \RuleSection{Results}
  \label{sec:results}

  \subsection{Recorded scan}
  \label{sec:real}

  \begin{figure}[htbp]
  \centering
  \begin{tikzpicture}
  \begin{axis}[width=0.62\textwidth, height=7.2cm, axis equal image,
      xmin=-1.6, xmax=2.0, ymin=-0.3, ymax=5.9,
      xlabel={lateral $z$ [m] (right $>0$)}, ylabel={forward $x$ [m]},
      label style={font=\small}, tick label style={font=\scriptsize},
      legend style={font=\scriptsize, at={(1.03,0.5)}, anchor=west, draw=none}]
    \addplot[only marks, mark=*, mark size=3pt, leftc] coordinates {(-0.562,2.849) (-0.827,4.882)};
    \addlegendentry{left cone}
    \addplot[only marks, mark=*, mark size=3pt, rightc] coordinates {(0.695,2.879) (0.457,3.977) (0.649,5.328)};
    \addlegendentry{right cone}
    \addplot[only marks, mark=o, mark size=3pt, black!50, thick] coordinates {(1.243,4.511)};
    \addlegendentry{unassigned}
    \addplot[very thick, orange!90!black] coordinates {(0,0) (-0.447,2.782) (-0.699,4.347)};
    \addlegendentry{racing line}
    \addplot[only marks, mark=triangle*, mark size=4pt, black] coordinates {(0,0)};
    \addlegendentry{car}
  \end{axis}
  \end{tikzpicture}
  \caption{Result on \code{cones.pcd} (top view). The sixth cone sits 0.6\,m outside the right
  boundary and fits neither side, so it is left unassigned rather than guessed.}
  \label{fig:real}
  \end{figure}

  Six cones are detected, with ICP residuals of 1.9--3.4\,cm RMS. Two are assigned to the left
  boundary and three to the right. With two cones on the left there are two gates. The optimal
  line through them is a straight line from the car (cost $\approx 0$), which is the expected
  minimum-curvature answer for such a short horizon.

  \subsection{Left/right split}

  \begin{table}[htbp]
  \small
  \centering
  \begin{tabular}{@{}lrr@{}}
  \toprule
  \textbf{Condition (10 tracks, 100 seeds each)} & \textbf{Previous version} & \textbf{Current} \\
  \midrule
  Clean tracks classified perfectly & 2 / 10 & \textbf{10 / 10} \\
  Accuracy on clean tracks & 7.7--100\,\% & \textbf{100\,\%} \\
  Perfect runs, 5\,cm position noise & 0\,\% & \textbf{100\,\%} \\
  Perfect runs, 15\,cm position noise & 0\,\% & 68--100\,\% \\
  Perfect runs, 10\,\% of cones missing & 0--79\,\% & 44--100\,\% \\
  \bottomrule
  \end{tabular}
  \caption{Left/right split on synthetic tracks. Ranges are over the 10 tracks; the weakest
  cases are the hairpin and the S-curve with missing cones.}
  \label{tab:lr}
  \end{table}

  When a decision is ambiguous, the current split leaves the cone unassigned rather than
  putting it on the wrong side. This keeps wrong boundaries out of the racing line.

  \subsection{Racing line}

  \begin{itemize}
    \item \textbf{Straight track:} the line is the centreline; maximum offset 0.000\,m.
    \item \textbf{Constant arc ($R=10$\,m):} the line hugs the outer edge. The end points sit exactly
      on the bound and interior points within 2.3\,\% of the half-width. That small gap is the
      true discrete optimum: it has a lower cost than putting every point on the bound.
    \item \textbf{Start at the car:} including the car in the cost reduces the initial turn from
      57.3$^\circ$ to 0.9$^\circ$.
    \item \textbf{All 10 synthetic tracks:} no bound violation, no reversal, forward progress at
      every gate; largest turn 24.6$^\circ$.
    \item \textbf{Drawn line:} it passes through every optimised point, in order, within $2\,\mu$m.
  \end{itemize}

  \subsection{Odometry}
  \label{sec:odometry}

  \begin{table}[htbp]
  \small
  \centering
  \begin{tabular}{@{}lrrrr@{}}
  \toprule
  & \multicolumn{2}{c}{\textbf{Unbounded (previous)}} & \multicolumn{2}{c}{\textbf{1\,m bound (current)}} \\
  \cmidrule(lr){2-3}\cmidrule(l){4-5}
  \textbf{True motion (forward, left, yaw)} & position err. & rotation err. & position err. & rotation err. \\
  \midrule
  0.5\,m, 0\,m, 0$^\circ$     & 0.837\,m & 10.07$^\circ$ & \textbf{0.002\,m} & \textbf{0.00$^\circ$} \\
  0.5\,m, +0.2\,m, +2$^\circ$ & 0.962\,m &  8.66$^\circ$ & \textbf{0.002\,m} & \textbf{0.00$^\circ$} \\
  0.3\,m, $-0.1$\,m, $-3^\circ$ & 0.873\,m & 10.13$^\circ$ & \textbf{0.002\,m} & \textbf{0.00$^\circ$} \\
  1.0\,m, 0\,m, +5$^\circ$    & 0.774\,m & 10.08$^\circ$ & \textbf{0.002\,m} & \textbf{0.00$^\circ$} \\
  \bottomrule
  \end{tabular}
  \caption{Odometry on synthetic scan pairs built from \code{first.pcd} with a known motion.
  With PCL's default unbounded correspondences, points with no counterpart in the other scan
  dominate the fit; with a 1\,m bound the motion is recovered.}
  \label{tab:odo}
  \end{table}

  On the recorded pair the current estimate is 0.41\,m forward and 0.27\,m left (0.49\,m in the
  ground plane), with a 6.3$^\circ$ rotation. ICP stops at the 70-iteration cap, and the residual
  over matched pairs is 0.46\,m RMS.

  \begin{warnbox}[Limitation: the recorded odometry pair cannot be verified]
  \small
  The two recorded scans are levelled with different hand-tuned pitch corrections
  ($-10^\circ$ and $-3^\circ$). Part of the estimated rotation and the 0.21\,m vertical motion
  most likely absorb that difference. The method is shown correct on synthetic pairs; the
  accuracy on this particular pair is unknown. Automatic levelling (ground-plane fit) and a
  motion prior would remove this uncertainty.
  \end{warnbox}

  \subsection{Cone detection}
  \label{sec:detection}

  On simulated scans (5 scenes, 3 seeds), the detector reaches \textbf{79\,\% precision and
  37\,\% recall}; the centroid error is 6.9\,cm, biased towards the sensor because only the
  near side of each cone is seen. Tracing every true cone through the stages shows three
  separate losses:
  \begin{enumerate}
    \item SOR with $k=50$ neighbours removes every cone beyond about 6\,m: such cones have fewer
      than 50 points, so their neighbourhood reaches the ground and they look like outliers.
    \item The 5\,cm clustering tolerance starts splitting cones at about 5\,m.
    \item The yaw gate rejects true cones. A cone is symmetric about its axis, so ICP cannot
      measure its yaw: a cone turned 30$^\circ$ about its axis is reported at $-7.3^\circ$. The
      gate accepts 85\,\% of simulated cone views, but also 95\,\% of cone-sized boxes and 100\,\%
      of buckets.
  \end{enumerate}

  The \emph{tilt} of the fitted axis is an observable quantity and separates the classes in
  simulation (cones 5.6--26.6$^\circ$, non-cones 30.7--56.8$^\circ$). Table~\ref{tab:variants}
  measures the effect of replacing the gate and retuning the two filters.

  \begin{table}[htbp]
  \small
  \centering
  \begin{tabular}{@{}lrrrrr@{}}
  \toprule
  \textbf{Detector variant} & \textbf{Precision} & \textbf{Recall} & \textbf{Sides ok} & \textbf{Line in track} & \textbf{Cones, recorded} \\
  \midrule
  Current (yaw $<13^\circ$) & 79\,\% & 37\,\% & 85\,\% & 59\,\% & 6 \\
  No yaw gate & 73\,\% & 52\,\% & 82\,\% & 80\,\% & 9 \\
  Tilt $<20^\circ$ & 94\,\% & 48\,\% & 98\,\% & 88\,\% & 7 \\
  Tilt $<20^\circ$, SOR $k=10$, clustering 8\,cm & 96\,\% & 70\,\% & 75\,\% & 88\,\% & 6 \\
  \bottomrule
  \end{tabular}
  \caption{Detector variants on simulated scans, and the number of cones each finds on the
  recorded scan. ``Line in track'' is the share of the racing line inside the true track,
  where a line was produced.}
  \label{tab:variants}
  \end{table}

  The tilt gate is clearly better in simulation. On the recorded scan, two candidates rejected
  by the current gate, at $(x,z)=(1.78,+0.55)$ and $(1.99,-0.65)$\,m, form a symmetric pair
  1.2\,m apart and about 1\,m before the first detected pair: very likely the nearest gate of
  cones, partly cut by the lower edge of the field of view. The tilt gate recovers one of them.
  Without labels this cannot be confirmed, so the production gate was left unchanged pending a
  visual check.

  \subsection{Runtime and resources}

  \begin{table}[htbp]
  \small
  \centering
  \begin{tabular}{@{}lrl@{}}
  \toprule
  \textbf{Stage} & \textbf{Time per scan} & \textbf{Notes} \\
  \midrule
  Cone detection (crop to ICP) & $221\pm5$\,ms & \\
  Obstacle extraction & $264\pm18$\,ms & second normal estimation and region growing \\
  Left/right split & 0.002\,ms & \\
  Racing line & 0.034\,ms & \\
  \textbf{Cone pipeline total} & $\bm{485\pm22}$\,\textbf{ms} & median 474\,ms; 2.1\,Hz; peak memory 120\,MiB \\
  \textbf{Odometry ICP} & \textbf{2.20\,s} & median of 10 runs; 0.45\,Hz; 92\,\% of one core; 108\,MiB \\
  \bottomrule
  \end{tabular}
  \caption{Apple M3, macOS 15.6, Release build (\code{-O3}), single thread. Cone pipeline: 10
  timed runs after a warm-up.}
  \label{tab:runtime}
  \end{table}

  % ======================================================================
  \RuleSection{Defects found by the validation}
  \label{sec:fixes}

  Three defects produced results that looked plausible on the recorded data but were wrong.
  Only tests with known answers exposed them.
  \begin{enumerate}
    \item \textbf{Left/right split on one coordinate.} A coordinate swap put the lateral position
      in a field the algorithm did not read, so distances and angles used the lateral offset
      alone. Fixed by working in the ground plane throughout; 2/10 $\to$ 10/10 clean tracks.
    \item \textbf{Racing line drawn differently from what was optimised.} Each drawn segment
      reused one offset for both ends, the last point was never drawn, and the optimiser
      ignored the car. The old curvature term also gave a U-turn zero cost. All four are fixed
      and covered by tests.
    \item \textbf{Odometry off by about 0.8\,m.} Unbounded correspondences, plus a sign inversion
      in the printed displacement. Fixed with a 1\,m bound and a single, tested displacement
      formula.
  \end{enumerate}

  % ======================================================================
  \RuleSection{Limitations and roadmap}
  \label{sec:roadmap}

  \begin{table}[htbp]
  \small
  \centering
  \renewcommand{\arraystretch}{1.1}
  \begin{tabularx}{\textwidth}{@{}L{0.8cm} L{5.0cm} X@{}}
  \toprule
  \textbf{Prio.} & \textbf{Item} & \textbf{Evidence and proposed action} \\
  \midrule
  1 & Cone recall at range & 37\,\% recall; SOR and clustering lose far cones. Make both
      range-aware, or downsample adaptively; target the 70\,\% measured with retuned filters. \\[.15cm]
  2 & Cone/non-cone decision & Yaw is not observable; replace it with a tilt gate after
      confirming the extra recorded cluster; add labelled scans for precision/recall by range. \\[.15cm]
  3 & Real-time throughput & 0.49\,s vs.\ 0.1\,s budget. Share normals and kd-trees between the
      two region-growing passes, downsample, parallelise. \\[.15cm]
  4 & Odometry on real data & Automatic levelling (ground-plane RANSAC \citep{fischler1981ransac}),
      motion prior, point-to-plane ICP or GICP \citep{segal2009gicp}; report convergence and
      inlier ratio. \\[.15cm]
  5 & Gate pairing & Gates pair the $i$-th left cone with the $i$-th right cone; a missed cone
      skews a gate. Pair by nearest opposite cone instead. \\[.15cm]
  6 & Code structure & Free functions with exported constants; move to classes with explicit
      parameters and ownership. \\
  \bottomrule
  \end{tabularx}
  \caption{Prioritised next steps, each tied to a measurement in this report.}
  \label{tab:roadmap}
  \end{table}

  % ======================================================================
  \RuleSection{Reproducing the results}
  \label{sec:repro}

  Dependencies: PCL 1.15, VTK 9.5, Eigen, NLopt 2.10, Boost and Qt 5 or 6 (GUI only); CMake
  $\ge$ 3.5 (4.0 for the odometry project). From each project folder:
  \begin{verbatim}
  cmake -DCMAKE_BUILD_TYPE=Release -G Ninja -S . -B build && cmake --build build
  ctest --test-dir build
  \end{verbatim}
  The two visual reports are written by:
  \begin{itemize}
    \item \code{eval\_left\_right -{}-pcd cones.pcd -{}-model ideal\_cone.ply} (left/right split);
    \item \code{eval\_pipeline -{}-model ideal\_cone.ply} (whole cone pipeline on simulated scans).
  \end{itemize}
  The repository README lists every command.

  The code footprint is 1,119 lines for the pipeline logic, 452 for the viewer, 239 for
  odometry, and 1,361 for tests and evaluation tools (non-blank lines).

  % ======================================================================
  \bibliographystyle{abbrvnat}
  \bibliography{references}

  \end{document}
